% chapters/11/section-01.tex
\section{Menghubungkan Dua Lapisan}
\label{sec:hubung-lapisan}

Bab~IX menghasilkan keputusan, dan Bab~X menghasilkan rute. Bagian ini
menghubungkan keduanya menjadi satu NPC yang utuh.

\subsection{Niat dan Rute}

Perhatikan apa yang dihasilkan masing-masing lapisan.

\begin{itemize}
  \item \textbf{Mesin keadaan pada Bab~IX} menghasilkan \textbf{niat}, yaitu
        jawaban atas pertanyaan ``apa yang sebaiknya saya lakukan?''. Jawabannya
        berupa keadaan, misalnya mengejar.
  \item \textbf{Pencarian jalur pada Bab~X} menghasilkan \textbf{rute}, yaitu
        jawaban atas pertanyaan ``bagaimana cara menuju ke sana?''. Jawabannya
        berupa rangkaian petak.
\end{itemize}

Pada Bab~IX, niat ``mengejar'' dilaksanakan dengan bergerak lurus ke arah
pemain. Pada bab ini, niat yang sama dilaksanakan dengan memanggil pencarian
jalur terlebih dahulu. Perubahan itu hanya menuntut satu tambahan pada fungsi
tindakan, sedangkan mesin keadaannya tidak berubah sama sekali.

\begin{kodekeluaran}
if self.state == "mengejar":
    self.selang_tersisa -= delta_seconds
    if self.selang_tersisa <= 0.0 or not self.jalur:
        self.hitung_jalur(petak_dari_piksel(pemain.position))
        self.selang_tersisa = self.selang
    self.ikuti_jalur(delta_seconds, CHASE_SPEED, dinding)
\end{kodekeluaran}

Susunan itu diperlihatkan pada Gambar~\ref{fig:arsitektur-integrasi}. Perhatikan
empat kotak pada gambar tersebut: mesin keadaan, pencarian jalur, simpanan
jalur, dan pengikut jalur. Keempatnya bekerja berurutan, dan masing-masing
memiliki tanggung jawab yang jelas.

\begin{figure}[htbp]
\centering
% chapters/11/figures/arsitektur-integrasi.tikz
% Hubungan mesin keadaan dan pencarian jalur di dalam satu NPC.
\begin{tikzpicture}[
  kotak/.style={
    rectangle, rounded corners=2pt, draw=ubPrimary, thick, fill=ubAlt,
    align=center, text width=36mm, inner sep=4pt, minimum height=11mm,
    font=\footnotesize
  },
  simpan/.style={kotak, draw=ubGreen, fill=ubGreenBg},
  panah/.style={-{Stealth[length=2.4mm]}, thick, draw=ubAccent},
  label/.style={font=\scriptsize, align=center}
]

% Lapisan keputusan.
\node[kotak] (fsm) at (0,0)
  {\textbf{Mesin keadaan}\\[2pt]\scriptsize berpatroli, mengejar, menyerang};

% Lapisan navigasi.
\node[kotak, below=16mm of fsm] (cari)
  {\textbf{Pencarian jalur A*}\\[2pt]\scriptsize dari petak sekarang ke petak pemain};

% Simpanan jalur dan pengikut jalur.
\node[simpan, right=14mm of cari] (jalur)
  {\textbf{Simpanan jalur}\\[2pt]\scriptsize rangkaian petak beserta\\waktu hitung ulang terakhir};
\node[kotak, above=16mm of jalur] (ikuti)
  {\textbf{Pengikut jalur}\\[2pt]\scriptsize bergerak menuju petak\\berikutnya pada jalur};

\draw[panah] (fsm) -- node[label, left, xshift=-2pt]
  {niat:\\kejar} (cari);
\draw[panah] (cari) -- node[label, below, yshift=-1pt]
  {jalur} (jalur);
\draw[panah] (jalur) -- node[label, left, xshift=-2pt]
  {dipakai} (ikuti);
\draw[panah] (ikuti) -- node[label, right, xshift=2pt]
  {gerak} (fsm);

\node[label, anchor=north, text width=140mm] at (3.4,-4.4)
  {Mesin keadaan menentukan niat, sedangkan pencarian jalur menentukan rute.
   Karena pencarian jalur mahal, hasilnya disimpan dan dipakai berkali-kali;
   pencarian diulang hanya menurut selang waktu tertentu atau ketika jalurnya
   sudah habis dilalui.};
\end{tikzpicture}

\caption{Hubungan mesin keadaan dan pencarian jalur di dalam satu NPC}
\label{fig:arsitektur-integrasi}
\end{figure}

\subsection{Kode Lengkap}

\kodepython{python/ch11_mini_proyek.py}{Mini proyek AI: mesin keadaan yang memanggil pencarian jalur}{lst:ch11-mini-proyek}

\subsection{Pencarian yang Tidak Dijalankan Setiap Bingkai}

Perhatikan syarat pada potongan kode di atas. Pencarian jalur dijalankan
apabila salah satu dari dua hal terjadi.

\begin{enumerate}
  \item \textbf{Selang waktu tertentu telah berlalu} sejak pencarian terakhir.
        Selang itu disebut selang hitung ulang, dan nilainya menjadi pokok
        Bagian~\ref{sec:batas-kinerja}.
  \item \textbf{Jalurnya sudah habis} karena seluruh petaknya telah dilalui.
        Keadaan itu terjadi ketika NPC telah sampai di tempat tujuan dan
        pemain belum berpindah, sehingga pencarian ulang menuntut perhitungan
        yang singkat saja.
\end{enumerate}

Syarat kedua itu mudah terlupakan, dan akibatnya NPC berhenti bergerak
walaupun pemainnya masih terlihat: jalurnya sudah habis dilalui, tetapi
pencarian ulang belum waktunya dijalankan.

\subsection{Menyimpan Hasil Pencarian}

Hasil pencarian disimpan pada dua tempat di dalam NPC, yaitu rangkaian
petaknya dan pewaktu selang hitung ulang.

\begin{kodekeluaran}
def hitung_jalur(self, tujuan_petak):
    hasil = pencarian.cari_jalur(pencarian.PETA, self.petak_sekarang(),
                                 tujuan_petak, "manhattan")
    self.jalur = hasil["jalur"][1:]     # petak pertama adalah petak sendiri
    self.statistik["hitung_ulang"] += 1
    self.statistik["petak_diperiksa"] += len(hasil["diperiksa"])
\end{kodekeluaran}

Perhatikan dua hal. Pertama, petak pertama pada hasil pencarian dibuang, sebab
petak itu adalah petak tempat NPC berada sekarang; menyimpannya membuat NPC
berusaha mencapai petaknya sendiri. Kedua, kedua penghitung statistik ditambah
di sini, dan itulah bahan penilaian pada Bagian~\ref{sec:menilai-mutu}.

\begin{catatabox}
Perhatikan baris \code{import ch10_astar as pencarian} pada awal program ini.
Pencarian jalur tidak ditulis ulang, melainkan dipakai dari berkas bab
sebelumnya. Kebiasaan memakai kembali kode yang sudah ada---bukan menyalinnya---
merupakan kebiasaan yang penting: perbaikan pada pencarian jalur cukup
dilakukan di satu tempat, dan seluruh pengguna pencarian itu ikut memperoleh
perbaikannya.
\end{catatabox}
